\documentclass[12pt,a4paper]{article}
\usepackage[utf8]{inputenc}
\usepackage[T1]{fontenc}
\usepackage[left=2.5cm,right=2.5cm,top=2cm,bottom=2cm]{geometry}
\usepackage{amsmath,amssymb}
\usepackage{enumitem}
\usepackage{extramarks}
\usepackage{fancyhdr}
\usepackage[ngerman]{babel}

\pagestyle{fancy}
\fancyhf{}
\fancyhead[L]{Programmierung WiSe 2024/2025}
\fancyhead[R]{Klausur - Haskell Vertiefung}
\fancyfoot[C]{\thepage}

\makeatletter
\def\cleardoublepage{\clearpage\if@twoside\ifodd\c@page\else%
\hbox{}%
\thispagestyle{empty}%
\newpage%
\if@twocolumn\hbox{}\newpage\fi\fi\fi}
\makeatother

\begin{document}

\begin{center}
    \textbf{\Large Klausur Programmierung - Haskell Vertiefung}\\
    \vspace{0.5cm}
    \textbf{Schwierigkeitsgrad: Schwer}\\
    \vspace{0.5cm}
    \textbf{Klausurdatum: März 2025}
\end{center}

\vspace{1cm}

\noindent
\textbf{Name:} \underline{\hspace{8cm}}\\
\textbf{Matrikelnummer:} \underline{\hspace{5cm}}\\
\textbf{Studiengang:} \underline{\hspace{5cm}}

\vspace{1cm}

\begin{center}
\begin{tabular}{|c|c|c|}
\hline
\textbf{Aufgabe} & \textbf{Punkte} & \textbf{Erreicht} \\ \hline
Aufgabe 1 & 15 & \\ \hline
Aufgabe 2 & 12 & \\ \hline
Aufgabe 3 & 18 & \\ \hline
Aufgabe 4 & 20 & \\ \hline
Aufgabe 5 & 15 & \\ \hline
Aufgabe 6 & 20 & \\ \hline
\textbf{Summe} & \textbf{100} & \\ \hline
\end{tabular}
\end{center}

\vspace{1cm}

\noindent
\textbf{Allgemeine Hinweise:}
\begin{itemize}
    \item Schreiben Sie mit dokumentenechten Stiften.
    \item Geben Sie für jede Aufgabe maximal eine Lösung an.
    \item Bei Täuschungsversuchen wird die Klausur mit 0 Punkten bewertet.
    \item Alle Hilfsmittel sind zugelassen.
\end{itemize}

\cleardoublepage

\section*{Aufgabe 1 (Typinferenz und Typklassen): (15 Punkte)}

Gegeben sei folgende Haskell-Programme:

\begin{verbatim}
class Monad m where
    (>>=) :: m a -> (a -> m b) -> m b
    return :: a -> m a

instance Monad Maybe where
    return = Just
    Nothing >>= _ = Nothing
    Just x >>= f = f x

instance Monad [] where
    return x = [x]
    xs >>= f = concat (map f xs)
\end{verbatim}

\begin{enumerate}
    \item \textbf{(8 Punkte)} Bestimmen Sie den allgemeinsten Typ für jede der folgenden Funktionen und begründen Sie Ihre Antwort:
    
    \begin{verbatim}
    f1 :: ?
    f1 x y = if x then [y] else []

    f2 :: ?
    f2 x = map ($ x)

    f3 :: ?
    f3 = do
        a <- [1,2,3]
        b <- [4,5]
        return (a,b)

    f4 :: ?
    f4 [] = Nothing
    f4 (x:xs) = Just (x, xs)
    \end{verbatim}

    \item \textbf{(7 Punkte)} Gegeben sei folgende Typdeklaration:
    \begin{verbatim}
    class Functor f => Applicative f where
        pure :: a -> f a
        (<*>) :: f (a -> b) -> f a -> f b
    \end{verbatim}
    
    Zeigen Sie, dass \texttt{Maybe} eine Instanz von \texttt{Applicative} ist, indem Sie \texttt{pure} und \texttt{<*>} für \texttt{Maybe} implementieren. Beweisen Sie anschließend, dass die \texttt{Applicative}-Gesetze erfüllt sind:
    \begin{enumerate}
        \item \texttt{pure id <*> v = v}
        \item \texttt{pure (.) <*> u <*> v <*> w = u <*> (v <*> w)}
        \item \texttt{pure f <*> pure x = pure (f x)}
        \item \texttt{u <*> pure y = pure (\$ y) <*> u}
    \end{enumerate}
\end{enumerate}

\cleardoublepage

\section*{Aufgabe 2 (Monaden und Monaden-Transformers): (12 Punkte)}

Gegeben sei folgende Datenstruktur für eine einfache Zustandsmonade:

\begin{verbatim}
newtype State s a = State { runState :: s -> (a, s) }

instance Functor (State s) where
    fmap f (State g) = State $ \s -> let (a, s') = g s in (f a, s')

instance Applicative (State s) where
    pure a = State $ \s -> (a, s)
    State f <*> State g = State $ \s ->
        let (h, s') = f s
            (a, s'') = g s'
        in (h a, s'')

instance Monad (State s) where
    State f >>= g = State $ \s ->
        let (a, s') = f s
            State h = g a
        in h s'
\end{verbatim}

\begin{enumerate}
    \item \textbf{(4 Punkte)} Implementieren Sie die folgenden Hilfsfunktionen:
    \begin{verbatim}
    get :: State s s
    put :: s -> State s ()
    modify :: (s -> s) -> State s ()
    \end{verbatim}
    
    \item \textbf{(8 Punkte)} Implementieren Sie einen \texttt{State}-Transformer namens \texttt{StateT}, der zusätzlich zu einem Zustand auch einen Fehlerwert verarbeiten kann:
    \begin{verbatim}
    newtype StateT s m a = StateT { runStateT :: s -> m (a, s) }
    \end{verbatim}
    
    Implementieren Sie die \texttt{Monad}-Instanz für \texttt{StateT s m}, wobei \texttt{m} eine Monad sein soll. Gehen Sie dabei wie folgt vor:
    \begin{enumerate}
        \item Definieren Sie die \texttt{Functor}-Instanz
        \item Definieren Sie die \texttt{Applicative}-Instanz
        \item Definieren Sie die \texttt{Monad}-Instanz unter Verwendung von \texttt{MonadTrans}:
        \begin{verbatim}
        class MonadTrans t where
            lift :: Monad m => m a -> t m a
        \end{verbatim}
    \end{enumerate}
\end{enumerate}

\cleardoublepage

\section*{Aufgabe 3 (Parser-Kombinatoren): (18 Punkte)}

In dieser Aufgabe entwickeln wir einen einfachen Parser-Kombinator. Ein Parser ist ein Programm, das eine Eingabe (einen String) analysiert und ein Ergebnis zurückgibt.

\begin{verbatim}
newtype Parser a = Parser { runParser :: String -> Maybe (a, String) }

instance Functor Parser where
    fmap f (Parser p) = Parser $ \input ->
        case p input of
            Nothing -> Nothing
            Just (a, rest) -> Just (f a, rest)

instance Applicative Parser where
    pure a = Parser $ \input -> Just (a, input)
    Parser pf <*> Parser px = Parser $ \input ->
        case pf input of
            Nothing -> Nothing
            Just (f, rest) ->
                case px rest of
                    Nothing -> Nothing
                    Just (x, rest') -> Just (f x, rest')

instance Monad Parser where
    Parser p >>= f = Parser $ \input ->
        case p input of
            Nothing -> Nothing
            Just (a, rest) -> runParser (f a) rest
\end{verbatim}

\begin{enumerate}
    \item \textbf{(5 Punkte)} Implementieren Sie folgende primitive Parser:
    \begin{verbatim}
    satisfy :: (Char -> Bool) -> Parser Char
    char :: Char -> Parser Char
    anyChar :: Parser Char
    \end{verbatim}
    
    \item \textbf{(5 Punkte)} Implementieren Sie folgende Kombinations-Parser:
    \begin{verbatim}
    -- Parsen von null oder mehr Vorkommen
    many :: Parser a -> Parser [a]
    
    -- Parsen von einem oder mehr Vorkommen
    many1 :: Parser a -> Parser [a]
    
    -- Alternative: erster erfolgreicher Parser
    (<|>) :: Parser a -> Parser a -> Parser a
    \end{verbatim}
    
    \item \textbf{(8 Punkte)} Implementieren Sie einen Parser für arithmetische Ausdrücke mit folgenden Regeln:
    \begin{itemize}
        \item \texttt{expr := term (('+' | '-') term)*}
        \item \texttt{term := factor (('*' | '/') factor)*}
        \item \texttt{factor := number | '(' expr ')'}
        \item \texttt{number := [0-9]+}
    \end{itemize}
    
    Der Parser soll eine Zahl als \texttt{Int} parsen. Implementieren Sie:
    \begin{verbatim}
    expr :: Parser Int
    term :: Parser Int
    factor :: Parser Int
    number :: Parser Int
    \end{verbatim}
    
    \textit{Hinweis:} Sie können \texttt{read} verwenden.
\end{enumerate}

\cleardoublepage

\section*{Aufgabe 4 (Laziness und unendliche Datenstrukturen): (20 Punkte)}

\begin{enumerate}
    \item \textbf{(6 Punkte)} Erklären Sie den Unterschied zwischen \emph{strikter} und \emph{nicht-strikter} Auswertung in Haskell. Geben Sie jeweils ein Beispiel, wo der Unterschied beobachtbar ist.
    
    \item \textbf{(7 Punkte)} Implementieren Sie die folgenden Funktionen unter Ausnutzung der laziness von Haskell:
    \begin{verbatim}
    -- Fibonacci-Zahlen als unendliche Liste
    fibs :: [Integer]
    
    -- Alle Primzahlen (Sieb des Eratosthenes)
    primes :: [Integer]
    
    -- Binomialkoeffizienten als Dreieck
    binom :: [[Integer]]
    \end{verbatim}
    
    \item \textbf{(7 Punkte)} Gegeben sei folgende Funktion:
    \begin{verbatim}
    data Stream a = Cons a (Stream a)
    
    headS :: Stream a -> a
    headS (Cons x _) = x
    
    tailS :: Stream a -> Stream a
    tailS (Cons _ xs) = xs
    
    takeS :: Int -> Stream a -> [a]
    takeS 0 _ = []
    takeS n (Cons x xs) = x : takeS (n-1) xs
    \end{verbatim}
    
    Implementieren Sie:
    \begin{verbatim}
    -- Euklidischer Algorithmus als Stream
    euclid :: Integer -> Integer -> Stream Integer
    
    -- Alle Teiler einer Zahl
    divisors :: Integer -> [Integer]
    
    -- Collatz-Funktion
    collatz :: Integer -> Integer
    collatz n
        | n == 1    = 1
        | even n    = n `div` 2
        | otherwise = 3 * n + 1
    
    -- Collatz-Sequenz ab n
    collatzSeq :: Integer -> [Integer]
    \end{verbatim}
\end{enumerate}

\cleardoublepage

\section*{Aufgabe 5 (Rekursion und funktionale Datenstrukturen): (15 Punkte)}

\begin{enumerate}
    \item \textbf{(8 Punkte)} Implementieren Sie einen binären Suchbaum in Haskell:
    \begin{verbatim}
    data Tree a = Empty
                | Node a (Tree a) (Tree a)
        deriving (Show, Eq)
    
    -- Einfügen eines Elements
    insert :: Ord a => a -> Tree a -> Tree a
    
    -- Suchen eines Elements
    contains :: Ord a => a -> Tree a -> Bool
    
    -- Bauminversion (Spiegelung)
    mirror :: Tree a -> Tree a
    
    -- Traversierungen
    inOrder :: Tree a -> [a]
    preOrder :: Tree a -> [a]
    postOrder :: Tree a -> [a]
    levelOrder :: Tree a -> [a]
    \end{verbatim}
    
    \item \textbf{(7 Punkte)} Implementieren Sie eine Rot-Schwarz-Baum-Datenstruktur mit den folgenden Operationen:
    \begin{verbatim}
    data Color = R | B
    
    data RBTree a = Empty
                  | Node Color a (RBTree a) (RBTree a)
    
    -- Balance-Transformationen
    balance :: RBTree a -> RBTree a
    
    -- Einfügen mit automatischer Balance
    insertRB :: Ord a => a -> RBTree a -> RBTree a
    \end{verbatim}
    
    Erklären Sie, welche der RB-Baum-Eigenschaften durch Ihre \texttt{balance}-Funktion erhalten bleiben.
\end{enumerate}

\cleardoublepage

\section*{Aufgabe 6 (Fortgeschrittene Konzepte): (20 Punkte)}

\begin{enumerate}
    \item \textbf{(6 Punkte)} Erklären Sie das Konzept der \emph{Monade} in Haskell. Was ist der Unterschied zwischen \texttt{Monad} und \texttt{MonadPlus}? Wann würde man \texttt{MonadPlus} verwenden?
    
    \item \textbf{(7 Punkte)} Gegeben sei folgende Monade-Definition:
    \begin{verbatim}
    newtype Writer w a = Writer { runWriter :: (a, w) }
    
    instance Monoid w => Monad (Writer w) where
        return a = Writer (a, mempty)
        Writer (a, w) >>= f = let Writer (b, w') = f a
                               in Writer (b, w `mappend` w')
    \end{verbatim}
    
    Implementieren Sie:
    \begin{verbatim}
    tell :: w -> Writer w ()
    listen :: Writer w a -> Writer w (a, w)
    pass :: Writer w (a, w -> w) -> Writer w a
    \end{verbatim}
    
    \item \textbf{(7 Punkte)} Implementieren Sie einen Parser, der sowohl Fehlermeldungen als auch Logs führen kann, unter Verwendung von \texttt{StateT} und einem geeigneten Error-Typ:
    \begin{verbatim}
    type ParseError = String
    
    type ParserM a = StateT String (WriterT [String] (Either ParseError)) a
    
    -- Führt den Parser aus
    runParserM :: ParserM a -> String -> Either ParseError (a, String, [String])
    
    -- Primitive Parser
    parserFail :: ParseError -> ParserM a
    parserGet :: ParserM Char
    parserPut :: String -> ParserM ()
    \end{verbatim}
\end{enumerate}

\vfill

\begin{center}
\textbf{--- Ende der Klausur ---}
\end{center}

\end{document}
